\begin{frame}{결정 경계를 눈으로: ``계단''에서 ``곡선''으로}
  \begin{columns}[T]
    \begin{column}{0.52\textwidth}
      \centering
      \includegraphics[width=\linewidth]{ch07_2_boundary_tree_vs_rf.png}
      \vskip0.3em
      {\footnotesize 왼쪽: 트리 하나(깊이 3), 오른쪽: 랜덤 포레스트(100개).
      2개 특징 300개 샘플.}
    \end{column}
    \begin{column}{0.45\textwidth}
      \begin{itemize}
        \item 왼쪽: \textbf{축에 평행한 직선 조각이 이어진 계단형} —
          ``$x_1 > t$?'' 같은 축에 평행한 질문만 하므로 필연적
        \item 오른쪽: 100개 트리의 \textbf{다수결} $\to$ 각 트리의 계단이
          조금씩 달라 \textbf{훨씬 부드러운 곡선}
        \item 이 ``부드러짐'' = 앞 절 \textbf{분산 감소의 시각화}:
          단일 트리의 날카로운 계단은 잡음에도 쉽게 모양이 변하지만
          (높은 분산), 100개 계단의 평균은 잡음이 \textbf{상쇄}되어
          안정적
      \end{itemize}
    \end{column}
  \end{columns}
  \vskip0.5em
  {\footnotesize ``트리가 하나면 계단, 숲이면 곡선''이라는 앙상블의 본질을
  2D 데이터에서 눈으로 직접 보는 장면.}
\end{frame}
